\documentclass[conference]{IEEEtran}
\IEEEoverridecommandlockouts
\usepackage{cite}
\usepackage{amsmath,amssymb,amsfonts}
\usepackage{graphicx}
\usepackage{textcomp}
\usepackage{xcolor}
\usepackage{booktabs}
\usepackage{url}

% ---- Drafting helper: red notes. Delete all \TODO{...} before submission.
\newcommand{\TODO}[1]{\textcolor{red}{[TODO: #1]}}

% Short version. The full version is kept in main_long.tex.

\begin{document}

\title{Fine-Tuning Small Language Models for Intent Classification of Vietnamese Customer Support Tickets}

\author{
\IEEEauthorblockN{\TODO{Member 1}}
\IEEEauthorblockA{\TODO{University}, Vietnam \\ \TODO{email}}
\and
\IEEEauthorblockN{\TODO{Member 2}}
\IEEEauthorblockA{\TODO{University}, Vietnam \\ \TODO{email}}
\and
\IEEEauthorblockN{\TODO{Member 3}}
\IEEEauthorblockA{\TODO{University}, Vietnam \\ \TODO{email}}
}

\maketitle

\begin{abstract}
Support tickets sent to Vietnamese super-apps are often written without
diacritics, with chat abbreviations or with English terms. We define 45
labels that combine 9 services and 26 issues, create 900 synthetic
tickets, and test each in six forms of clean and non-standard text.
PhoBERT-base and Qwen2.5 models with 1.5 and 3 billion parameters are
fine-tuned on clean text only. The 3-billion-parameter model reaches 90.1
percent macro-F1 on clean tickets against 81.0 percent for PhoBERT, and its
gains are significant on every form. Without diacritics, PhoBERT keeps
about one fifth of its clean score and the larger model nearly two thirds.
\end{abstract}

\begin{IEEEkeywords}
intent classification, small language models, LoRA, Vietnamese, noisy text
\end{IEEEkeywords}

\section{Introduction}
\label{sec:intro}
Vietnamese super-apps handle rides, food, groceries, parcels, e-wallets
and memberships in one app, so one support channel must route tickets about
many services to the right team \cite{fodor2026beyond}. Tickets are often
non-standard (missing diacritics, teencode, Telex typing errors, English
terms), public Vietnamese datasets target chatbots
\cite{tran2020understanding}, and API models send personal data to outside
providers \TODO{cite Law No.~91/2025/QH15}. We ask whether small language
models fine-tuned with LoRA beat a fine-tuned Vietnamese encoder (RQ1), how
much scores drop on non-standard text when training uses clean text only
(RQ2), and which errors remain (RQ3).

\section{Related Work}
\label{sec:related}
Vietnamese intent detection has mostly targeted chatbots and flight
booking \cite{tran2020understanding,nguyen2022intent,nguyen2022advanced},
usually with PhoBERT \cite{nguyen2020phobert}. Customer-service studies
classify e-mails, chats and tickets
\cite{borg2021email,bonechi2024enhancing,pawlik2022method,rizou2023efficient,nowakbrzezinska2025hybrid,rebelo2022automatic,ferrera2025linguistics},
and LLMs have been used to create training data or compared with smaller
classifiers \cite{benayas2024enhancing,madrueno2025novel,alexander2025few,ouaddi2025assessing,pai2026finetuning}.
None of them compares small models tuned with LoRA
\cite{hu2022lora,qwen2024qwen25} with PhoBERT on noisy Vietnamese tickets.

\section{Dataset}
\label{sec:dataset}
From public help pages of Vietnamese super-apps we compiled 54 support
intents and split each into a service (9 in total) and an issue (26 in
total). Not every issue applies to every service, so only 45 of the
$9\times26$ pairs are valid labels, e.g., \texttt{food/charged\_but\_failed}.
An LLM (Claude Opus 5.5) wrote 20 tickets per label:
$45\times20=900$ tickets of 12.2
words on average, without duplicates; 36 do not name the service. For each
label, 14 tickets go to training, 3 to validation and 3 to test, giving
630/135/135 tickets. Each test ticket is evaluated in six rule-based
variants: clean, no diacritics, teencode (chat abbreviations), Telex typos
(raw keystrokes in one or two words), English terms (support words replaced
by English; this changes 345 of 900 tickets) and mixed (the last three
combined).

\section{Method}
\label{sec:method}
\subsection{Models}
\textbf{PhoBERT-base-v2} \cite{nguyen2020phobert} reads the word-segmented
ticket $x$ (segmented with underthesea) and turns it into one vector; all its weights
are trained. The \emph{flat} head turns this vector into 45 probabilities
and picks the highest. The \emph{hierarchical} head gives two sets of
probabilities, $p(s\mid x)$ over the 9 services and $p(i\mid x)$ over the 26
issues, and is trained with the sum of the two cross-entropy losses. It
predicts the valid pair with the highest combined score:
\begin{equation}
\hat{y}=\arg\max_{(s,i)\in V}\big[\log p(s\mid x)+\log p(i\mid x)\big],
\label{eq:hier}
\end{equation}
where $V$ is the set of 45 valid pairs. Adding the logarithms is the same
as multiplying the two probabilities.

\textbf{Qwen2.5-1.5B- and 3B-Instruct} \cite{qwen2024qwen25} read a short
instruction and the ticket and write the label as text. With LoRA
\cite{hu2022lora}, each weight matrix $W$ of size $d\times k$ stays fixed
and only a small correction is learned:
\begin{equation}
W' = W + \frac{\alpha}{r}\,BA,\qquad B\in\mathbb{R}^{d\times r},\;A\in\mathbb{R}^{r\times k},
\label{eq:lora}
\end{equation}
with $r=16$ and $\alpha=32$. Each matrix therefore gains $r(d+k)$ trainable
values; for example, a $1536\times1536$ attention matrix of the 1.5B model
gains $16\times3072=49{,}152$. Adding this over the seven matrices of every
layer gives 18.5M trainable values for the 1.5B model and 29.9M for the 3B
model, about 1\% of each model. Training minimizes
\begin{equation}
\mathcal{L}=-\sum_{t}\log p(y_t\mid x,y_{<t})
\label{eq:loss}
\end{equation}
over the label tokens $y_t$ only, so the model learns to answer rather
than to copy the ticket. When predicting, the model takes the most likely
token at each step, but only among tokens that continue one of the 45
labels, so every output is a valid label.

\subsection{Training}
PhoBERT uses a learning rate of $2\times10^{-5}$ for 10 epochs. Qwen uses
$2\times10^{-4}$ for 5 epochs, a larger rate because LoRA trains few values
that start from zero; it runs in bf16 on one T4 GPU. Both use 16 tickets
per update; for Qwen this is two batches of 8 added together, because of
GPU memory. These are common default values and were not tuned. After each
epoch we compute macro-F1 on the validation set and keep the best epoch;
the test set is never used for any choice. All models are trained on clean
text only. Each setting is run three times (seeds 42, 43, 44), and we
report the mean of the three runs.

\subsection{Evaluation}
For each label $c$,
\begin{equation}
F1_c=\frac{2\,TP_c}{2\,TP_c+FP_c+FN_c},\qquad
\text{Macro-F1}=\frac{1}{45}\sum_{c=1}^{45}F1_c,
\label{eq:f1}
\end{equation}
where $TP_c$ counts tickets of $c$ predicted as $c$, $FP_c$ tickets of other
labels predicted as $c$, and $FN_c$ tickets of $c$ predicted as something
else. Every label counts equally, whatever its size. Service F1 and issue F1
apply (\ref{eq:f1}) to the part before or after the slash.

The test set has only 135 tickets, so a difference between two models may
be due to chance. We check this with a paired bootstrap
\cite{efron1993bootstrap}: draw 135 test tickets with replacement, compute
the macro-F1 of both models on the same draw and take the difference;
repeat 2{,}000 times; the 2.5th and 97.5th percentiles of the 2{,}000
differences form a 95\% interval. If the interval does not contain 0, we
call the difference significant.

\section{Results}
\label{sec:results}
Table~\ref{tab:main} gives macro-F1 for each test variant; ``Noisy'' is the
mean of the five non-clean columns, e.g.,
$(56.3+88.0+82.3+90.0+61.3)/5=75.6$ for Qwen 3B.

\begin{table}[t]
\caption{Macro-F1 (\%) by Test Variant, Mean of Three Runs}
\label{tab:main}
\centering
\footnotesize
\setlength{\tabcolsep}{2.5pt}
\begin{tabular}{lccccccc}
\toprule
\textbf{Model} & \textbf{Clean} & \textbf{No diac.} & \textbf{Teen.} & \textbf{Telex} & \textbf{English} & \textbf{Mixed} & \textbf{Noisy} \\
\midrule
PhoBERT flat  & 79.1 & 16.0 & 78.2 & 63.2 & 74.2 & 19.3 & 50.2 \\
PhoBERT hier. & 81.0 & 16.1 & 78.6 & 64.2 & 74.6 & 16.6 & 50.0 \\
Qwen 1.5B     & 87.0 & 34.7 & 86.0 & 73.3 & 87.8 & 32.1 & 62.8 \\
Qwen 3B       & \textbf{90.1} & \textbf{56.3} & \textbf{88.0} & \textbf{82.3} & \textbf{90.0} & \textbf{61.3} & \textbf{75.6} \\
\bottomrule
\end{tabular}
\end{table}

\textbf{Clean text.} Qwen 3B is $90.1-81.0=9.1$ points above hierarchical
PhoBERT (95\% interval $[4.6, 16.1]$) and Qwen 1.5B is 6.1 points above
($[1.2, 13.0]$); both intervals exclude 0. Flat and hierarchical PhoBERT
differ by 1.9 points, but their interval $[-7.6, 3.2]$ contains 0, so
neither head is clearly better. Qwen 3B also raises service F1 from 91.5 to
94.6 and issue F1 from 84.4 to 93.8.

\textbf{Non-standard text.} Teencode lowers every model by at most 2.4
points (hierarchical PhoBERT: $81.0\rightarrow78.6$). English terms lower
PhoBERT by 4.9--6.4 points and Qwen by less than 1. Telex typos lower
PhoBERT by 16--17 points ($81.0\rightarrow64.2$) and Qwen 3B by 7.8
($90.1\rightarrow82.3$). Missing diacritics hurt most: PhoBERT keeps
$16.1/81.0=20\%$ of its clean score, Qwen 1.5B $34.7/87.0=40\%$ and Qwen 3B
$56.3/90.1=62\%$. The Qwen gains over PhoBERT are significant on all six
variants. A likely reason is that PhoBERT splits words using diacritics and
was trained on accented text, while Qwen saw unaccented Vietnamese during
pre-training.

\textbf{Errors.} We sort each prediction into four groups: correct, wrong
issue only, wrong service only, or both wrong (Fig.~\ref{fig:errors}).
Without diacritics, 41\% of hierarchical PhoBERT predictions are wrong on
both parts, against 13\% for Qwen 3B. Six test tickets are wrong for all
four models; four of them do not name the service, e.g., ``I was charged
although the transaction failed'' (gold \texttt{mart}). The test set has
five such tickets, i.e., 15 predictions over three runs: Qwen 3B finds the
right service in 2 (13\%), PhoBERT in none.

\begin{figure}[t]
\centering
\includegraphics[width=\columnwidth]{fig_error_types.pdf}
\caption{Share of predictions in each error group on clean and
no-diacritics tickets, over all three runs.}
\label{fig:errors}
\end{figure}

\section{Discussion}
\label{sec:discussion}
Fine-tuned small language models beat PhoBERT (RQ1). Their scores drop much
less without diacritics, and the larger model drops least (RQ2). Most
remaining errors come from tickets that do not name the service; these need
conversation context or a human agent (RQ3). Limitations: tickets are
synthetic and not yet checked by people; noise comes from rules; the test
set has three tickets per label; PhoBERT and Qwen 1.5B chose epochs near the
end of training, so they may need more epochs; speed and prompt-only LLMs
were not measured.

\section{Conclusion}
\label{sec:conclusion}
On a new two-level dataset of Vietnamese super-app tickets, Qwen2.5-3B with
LoRA reached 90.1\% macro-F1 on clean text and 75.6\% on average over the
five noisy variants, against 81.0\% and 50.0\% for PhoBERT. Future work
includes real tickets, training with noisy text, speed measurement and
multi-turn context.

\bibliographystyle{IEEEtran}
\bibliography{references}

\end{document}
